
\documentclass[12pt, a4paper]{article}

\usepackage[margin=1.2in]{geometry}
\usepackage{booktabs}
\usepackage{hyperref}
\usepackage{enumitem}
\usepackage{parskip}
\usepackage{setspace}
\usepackage{amsmath}
\usepackage{amssymb}
\usepackage{xcolor}
\usepackage{natbib}
\usepackage{titlesec}
\usepackage{fancyhdr}
\usepackage{listings}
\usepackage{graphicx}

\onehalfspacing

\hypersetup{colorlinks=true, linkcolor=black, citecolor=black, urlcolor=blue}

\pagestyle{fancy}
\fancyhf{}
\rhead{\small Bank Fragility Replication}
\lhead{\small \today}
\cfoot{\thepage}
\renewcommand{\headrulewidth}{0.4pt}

\titleformat{\section}{\large\bfseries}{\thesection.}{0.5em}{}
\titleformat{\subsection}{\normalsize\bfseries}{\thesubsection}{0.5em}{}

\definecolor{codegray}{gray}{0.95}
\lstset{
    backgroundcolor=\color{codegray},
    basicstyle=\ttfamily\footnotesize,
    breaklines=true,
    frame=single,
    framerule=0pt,
}

%% ─────────────────────────────────────────────────────────────────────────────
\begin{document}

\begin{titlepage}
    \centering
    \vspace*{1.5cm}

    {\LARGE\bfseries Replicating U.S. Bank Fragility:\\[0.3em]
    Mark-to-Market Losses and Uninsured Depositor Runs \par}

    \vspace{0.5cm}
    {\large An Open-Source Pipeline for Updated Estimation \par}

    \vspace{1.5cm}
    {\normalsize\itshape
    Replication of Jiang, Matvos, Piskorski, and Seru (2023),\\
    ``Monetary Tightening and U.S. Bank Fragility in 2023:\\
    Mark-to-Market Losses and Uninsured Depositor Runs?''\\[0.3em]
    \textit{Journal of Political Economy}, NBER Working Paper 31048
    \par}

    \vspace{2cm}
    {\normalsize
    \textbf{Report Date:} 12/31/2025 \\[0.3em]
    \textbf{Shock Window:} 2020-01-01 to 2025-12-31 \\[0.3em]
    \textbf{RMBS Multiplier:} 1.25
    \par}

    \vfill
    {\small
    \textit{Data sources: FFIEC Call Reports (public), FRED API (St.\ Louis Fed),
    Yahoo Finance.\\
    No WRDS or proprietary subscription required.}
    \par}
\end{titlepage}

%% ─────────────────────────────────────────────────────────────────────────────
\begin{abstract}
\noindent
We replicate the bank fragility analysis of Jiang, Matvos, Piskorski, and Seru (2023),
who document that the 2022--2023 monetary tightening cycle caused substantial
mark-to-market losses across the U.S. banking system and left a significant share
of banks vulnerable to uninsured depositor runs. Our contribution is twofold.
First, we reconstruct the full analytical pipeline using exclusively open-source
data --- the FFIEC public bulk download portal, the Federal Reserve's FRED API,
and Yahoo Finance --- replacing the proprietary WRDS database used in the original
paper. Second, we package this pipeline as a fully automated, reproducible workflow
that any researcher can execute for any Call Report date without manual data
collection. The methodology follows the original paper: bank assets are mapped to
maturity buckets, market shocks are derived from actual Treasury yield movements
over a user-specified window, and mark-to-market losses are estimated via duration
approximation. The primary fragility measure --- uninsured deposits as a share of
mark-to-market adjusted assets --- is computed for every U.S. commercial bank and
reported separately for small banks, large non-GSIB banks, and
Globally Systemically Important Banks (GSIBs).
\end{abstract}

\newpage
\tableofcontents
\newpage

%% ─────────────────────────────────────────────────────────────────────────────
\section{Introduction}

The collapse of Silicon Valley Bank (SVB) in March 2023 renewed academic and
policy interest in the interest rate risk embedded in bank balance sheets.
Jiang, Matvos, Piskorski, and Seru (2023) --- hereafter JMPS --- provided the
first comprehensive quantification of this risk across the entire U.S. banking
system. Their central finding was striking: after marking bank assets to market
using prevailing interest rates, aggregate bank assets fell by approximately
\$2.2 trillion, and roughly 190 banks with \$300 billion in assets became
technically insolvent on a mark-to-market basis. More critically, a large
fraction of these banks held substantial uninsured deposits, making them
vulnerable to self-fulfilling depositor runs even if their accounting book
values remained positive.

The JMPS paper relies on the Wharton Research Data Services (WRDS) platform,
a commercial database that requires an expensive institutional subscription.
This creates a significant barrier to replication and extension, particularly
for researchers at institutions without WRDS access and for practitioners
wishing to update the analysis as new Call Report quarters become available.

\textbf{Our contribution} is to reconstruct the JMPS methodology entirely
from publicly available data sources:

\begin{itemize}
    \item \textbf{FFIEC Call Reports} --- The Federal Financial Institutions
    Examination Council publishes quarterly Call Report data for all U.S.
    commercial banks free of charge through its bulk download portal.
    This is the same underlying data used in WRDS.

    \item \textbf{FRED API} --- The St.\ Louis Federal Reserve's FRED database
    provides the full history of Treasury constant maturity yields free of charge
    via a public API, replacing any proprietary yield data.

    \item \textbf{Yahoo Finance} --- MBS market conditions are proxied using
    publicly traded ETFs (SPMB for residential MBS, CMBS for commercial MBS),
    replacing the private MBS pricing data in the original paper.
\end{itemize}

Beyond data accessibility, we package the entire pipeline as an automated
workflow using the \texttt{doit} task runner. The pipeline handles dependency
tracking: it re-runs only the steps whose inputs have changed since the last
execution. This means a researcher can update the analysis to a new Call Report
quarter by changing a single configuration variable and running one command.

The remainder of this paper is structured as follows. Section~\ref{sec:data}
describes the data sources and construction. Section~\ref{sec:method} presents
the methodology for computing mark-to-market losses and the fragility measure.
Section~\ref{sec:analysis} reports the key results. Section~\ref{sec:pipeline}
describes the computational pipeline. Section~\ref{sec:conclusion} concludes.

%% ─────────────────────────────────────────────────────────────────────────────
\section{Data}
\label{sec:data}

\subsection{FFIEC Call Reports}

The primary data source is the FFIEC Call Report, a regulatory filing submitted
quarterly by every U.S. commercial bank to its federal supervisor. The bulk download
portal at \url{https://cdr.ffiec.gov} provides ZIP archives containing all schedules
for a given reporting period in tab-delimited format.

For this replication we use the Call Report as of 12/31/2025.
The following schedules are extracted from the ZIP archive:

\begin{center}
\begin{tabular}{lp{9cm}}
\toprule
Schedule & Content \\
\midrule
RC   & Consolidated Balance Sheet: total assets, total liabilities, equity \\
RCA  & Supplementary balance sheet items \\
RCB  & Securities: detailed breakdown by type and maturity \\
RCCI & Regulatory capital \\
RCE  & Equity capital components \\
\bottomrule
\end{tabular}
\end{center}

Each schedule contains thousands of MDRM-coded variables. Banks are identified
by their RSSD ID, a unique identifier assigned by the Federal Reserve. The
sample covers all U.S. commercial banks that filed a Call Report for the relevant
quarter.

\subsection{Bank Classification}

Following JMPS, we classify banks into three groups based on total assets and
systemic importance:

\begin{itemize}
    \item \textbf{Small banks:} Total assets below \$1.384 billion. These banks
    are below the regulatory threshold that triggers enhanced prudential standards
    under the post-Dodd-Frank framework.
    \item \textbf{Large non-GSIB banks:} Total assets at or above \$1.384 billion,
    excluding GSIBs. These banks face enhanced supervision but are not subject to
    the most stringent GSIB requirements.
    \item \textbf{GSIBs (Globally Systemically Important Banks):} The 37 institutions
    designated by the Financial Stability Board as systemically important. These are
    identified by their RSSD IDs and include the eight U.S. bank holding companies
    on the annual GSIB list.
\end{itemize}

\subsection{Treasury Yields}

Interest rate shocks are computed from daily Treasury constant maturity yields
obtained from the FRED database maintained by the Federal Reserve Bank of St.\
Louis. We use six maturities to span the yield curve: 1-year (DGS1), 3-year
(DGS3), 5-year (DGS5), 10-year (DGS10), 20-year (DGS20), and 30-year (DGS30).
The full history is available from 1962 onwards, with the 20-year series having
a gap from 1987 to 1993 when the Treasury suspended issuance of 20-year bonds.

\subsection{MBS Market Conditions}

The original JMPS paper uses private MBS pricing data to calibrate the additional
haircut applied to mortgage-backed securities. We proxy MBS market conditions
using two exchange-traded funds:

\begin{itemize}
    \item \textbf{SPMB} (SPDR Portfolio Mortgage Backed Bond ETF): tracks the
    Bloomberg U.S. MBS Index, providing exposure to agency residential
    mortgage-backed securities.
    \item \textbf{CMBS} (iShares CMBS ETF): tracks commercial mortgage-backed
    securities issued or guaranteed by U.S. government-sponsored enterprises.
\end{itemize}

Daily adjusted close prices are downloaded from Yahoo Finance. An RMBS multiplier
of 1.25 is applied to account for convexity risk --- the
tendency of MBS to extend in duration when interest rates rise as prepayment
speeds decline.

%% ─────────────────────────────────────────────────────────────────────────────
\section{Methodology}
\label{sec:method}

\subsection{Asset Maturity Bucketing}

A central challenge in estimating mark-to-market losses is that Call Reports do
not report the full maturity distribution of all asset categories. JMPS address
this by combining directly reported maturity breakdowns (available for securities)
with stylized maturity distributions for loan categories.

We follow the same approach. Bank assets are mapped to six maturity buckets:

\begin{center}
\begin{tabular}{llr}
\toprule
Bucket & Maturity Range & Duration Midpoint \\
\midrule
\texttt{lt1y}    & $< 1$ year      & 0.5 years  \\
\texttt{1\_3y}   & 1--3 years      & 2.0 years  \\
\texttt{3\_5y}   & 3--5 years      & 4.0 years  \\
\texttt{5\_10y}  & 5--10 years     & 7.5 years  \\
\texttt{10\_15y} & 10--15 years    & 12.5 years \\
\texttt{15plus}  & $> 15$ years    & 22.0 years \\
\bottomrule
\end{tabular}
\end{center}

For RMBS, Call Report schedule RCB provides actual holdings by maturity bucket,
which are used directly. For Treasury securities, CMBS, ABS, and other securities,
stylized maturity weights are applied to distribute total holdings across buckets.
For loan categories, similar stylized weights calibrated to observed loan maturity
distributions are applied.

\subsection{Market Shock Derivation}

Rather than using a hypothetical shock scenario, we derive the interest rate shock
empirically from actual Treasury yield movements over the shock window
(2020-01-01 to 2025-12-31).

For each maturity bucket, we interpolate a representative yield at the bucket
midpoint from the six available FRED maturities using linear interpolation:

\[
    y(\tau) = y_i + \frac{\tau - \tau_i}{\tau_{i+1} - \tau_i}(y_{i+1} - y_i)
    \quad \text{for } \tau_i \leq \tau \leq \tau_{i+1}
\]

where $\tau$ is the target maturity (bucket midpoint), and $\tau_i$, $\tau_{i+1}$
are the adjacent observed FRED maturities with yields $y_i$, $y_{i+1}$.

The shock for each bucket is then:
\[
    \Delta y_k = y_k(\text{end date}) - y_k(\text{start date})
\]

This approach ensures the shocks reflect actual market conditions over the chosen
period rather than a hypothetical scenario. Changing \texttt{MARKET\_START\_DATE}
and \texttt{MARKET\_END\_DATE} in the configuration allows the researcher to study
any historical episode.

\subsection{Mark-to-Market Loss Estimation}

Following JMPS, mark-to-market losses are estimated using the standard
duration approximation. For a given asset held in bucket $k$:

\[
    \text{Loss}_{i,k} = H_{i,k} \times \Delta y_k \times D_k \times m_k
\]

where:
\begin{itemize}[noitemsep]
    \item $H_{i,k}$ is bank $i$'s holdings in maturity bucket $k$ (in thousands of dollars)
    \item $\Delta y_k$ is the yield shock for bucket $k$ (in decimal form)
    \item $D_k$ is the duration midpoint of bucket $k$ (in years)
    \item $m_k$ is the asset-class multiplier: $m_k = 1.25$ for RMBS
    and residential mortgages (to account for convexity risk), and $m_k = 1$ otherwise
\end{itemize}

Total mark-to-market loss for bank $i$ is:
\[
    L_i = \sum_{k} \left(
        L^{\text{RMBS}}_{i,k} + L^{\text{res.mtg}}_{i,k} + L^{\text{tsy}}_{i,k}
        + L^{\text{other assets}}_{i,k} + L^{\text{other loans}}_{i,k}
    \right)
\]

Mark-to-market adjusted assets are:
\[
    A^{mm}_i = A_i - L_i
\]

where $A_i$ is total book-value assets of bank $i$.

\subsection{Fragility Measure}

The primary fragility measure, following JMPS, is the ratio of uninsured deposits
to mark-to-market adjusted assets:

\[
    \text{Fragility}_i = \frac{\text{Uninsured Deposits}_i}{A^{mm}_i}
\]

This ratio captures the share of potentially run-prone liabilities relative to the
bank's true economic net worth. A ratio above 1 implies that even a partial run
by uninsured depositors could exhaust the bank's assets. A ratio approaching 1
signals meaningful vulnerability.

Uninsured deposits are defined as total domestic deposits minus the FDIC-insured
portion. For the insured portion, we sum the relevant Call Report variables
(RCONHK05, RCONMT91, RCONMT87) which capture estimated insured deposits.

%% ─────────────────────────────────────────────────────────────────────────────
\section{Analysis and Results}
\label{sec:analysis}

\subsection{Aggregate Mark-to-Market Losses}

The mark-to-market loss for the U.S. banking system over the shock window
(2020-01-01 to 2025-12-31) is computed by
summing individual bank losses across all institutions. The magnitude of these
losses reflects the combination of (i) the size of the rate shock at each maturity,
(ii) the duration of bank asset portfolios, and (iii) the concentration of holdings
in long-duration instruments such as long-term Treasuries and RMBS.

Table~1 reports the full set of fragility statistics broken down by bank category.
The columns correspond to all banks, small banks (assets $<$ \$1.384B), large
non-GSIB banks, and GSIBs. For each group we report:

\begin{itemize}[noitemsep]
    \item The aggregate dollar loss summed across all banks in the group
    \item The mean and standard deviation of bank-level losses in dollars
    \item Portfolio composition: shares of RMBS, Treasury and other securities,
    residential mortgages, and other loans in total interest-rate-sensitive assets
    \item Loss/Asset: mean mark-to-market loss as a percentage of total assets
    \item Uninsured Deposit/MM Asset: the primary fragility measure
    \item Number of banks
\end{itemize}

\subsection{Cross-Sectional Heterogeneity}

A key finding of JMPS is that fragility is not uniformly distributed. Small banks
tend to hold a higher share of their assets in shorter-duration instruments
(short-term loans, adjustable-rate mortgages) and are therefore less exposed to
the rate shock on a per-dollar-of-assets basis. However, they also tend to have
a higher share of uninsured deposits relative to their size in some cases.

Large non-GSIB banks --- the category that includes institutions like SVB ---
are most heterogeneous. Some concentrated heavily in long-duration securities
(particularly agency RMBS and long-term Treasuries) funded by a small, concentrated
depositor base with a high fraction of uninsured deposits, making them acutely
vulnerable. Others maintained more diversified portfolios.

GSIBs, despite their size, benefit from their diversified global operations,
larger equity cushions, and the implicit government support embedded in their
designation, which reduces depositor run risk even when mark-to-market losses
are substantial.

\subsection{The Role of Asset Composition}

The portfolio composition columns in Table~1 reveal why certain bank groups
are more exposed. RMBS receive an additional multiplier of 1.25
relative to Treasuries of the same duration because rising rates simultaneously
reduce RMBS prices and extend their effective duration (as refinancing activity
drops). Banks with high RMBS concentrations therefore suffer disproportionately
larger losses for the same rate shock.

Long-duration Treasuries (in the \texttt{10\_15y} and \texttt{15plus} buckets)
are the other major source of losses. The 10-year Treasury yield rose by
approximately 300 basis points between early 2022 and late 2023, implying a
mark-to-market loss of roughly 15--20\% on 10-year bonds --- larger than the
equity capital buffers held by many community banks.

\subsection{Sensitivity to the Shock Window}

One advantage of our approach over a fixed hypothetical shock is that the results
are directly tied to observed market movements. The shock window can be set to any
historical episode of interest:

\begin{itemize}[noitemsep]
    \item \textbf{2022--2023 Fed hiking cycle} (\texttt{MARKET\_START\_DATE=2022-01-01},
    \texttt{MARKET\_END\_DATE=2023-12-31}): replicates the JMPS scenario most closely.
    \item \textbf{Post-COVID normalisation} (\texttt{MARKET\_START\_DATE=2020-01-01},
    \texttt{MARKET\_END\_DATE=2023-12-31}): captures the full cycle from near-zero
    rates through the tightening episode.
    \item \textbf{2018 hiking cycle}: tests bank resilience to an earlier, smaller
    tightening episode.
\end{itemize}

Changing these dates in \texttt{.env} and re-running \texttt{doit} regenerates
all results automatically, with only the affected downstream tasks re-executed.

%% ─────────────────────────────────────────────────────────────────────────────
\section{Computational Pipeline}
\label{sec:pipeline}

\subsection{Design Principles}

The pipeline is designed around three principles: (i) full reproducibility ---
every output is deterministically generated from the raw inputs; (ii) efficiency
--- tasks are skipped if their outputs are already up to date relative to their
inputs; and (iii) accessibility --- no proprietary data subscription is required.

The task runner is \texttt{doit}, a Python-based build system analogous to GNU
\texttt{make}. Each script is defined as a task with explicit \texttt{file\_dep}
(inputs) and \texttt{targets} (outputs). \texttt{doit} tracks file modification
times and re-executes only tasks whose dependencies have changed.

\subsection{Task Dependency Graph}

\begin{center}
\begin{tabular}{lll}
\toprule
Task & Inputs & Outputs \\
\midrule
\texttt{pull\_ffiec}          & \texttt{.env}                          & FFIEC ZIP \\
\texttt{pull\_gsib}           & \texttt{.env}                          & \texttt{gsib\_list.parquet} \\
\texttt{pull\_treasury}       & \texttt{.env}                          & \texttt{treasury\_yields.parquet} \\
\texttt{pull\_mbs}            & \texttt{.env}                          & \texttt{mbs\_etfs.parquet} \\
\texttt{compute\_yield\_shocks} & Treasury parquet, \texttt{.env}      & \texttt{market\_shocks.parquet} \\
\texttt{process\_ffiec}       & FFIEC ZIP, GSIB list, \texttt{.env}    & Bank panel, summary stats \\
\texttt{make\_table\_1}       & Bank panel, shocks, \texttt{.env}      & Table 1 CSV/\TeX{} \\
\texttt{convert\_notebook}    & \texttt{eda.ipynb}                     & \texttt{eda.html} \\
\texttt{generate\_pdf}        & This script, \texttt{.env}             & This PDF \\
\bottomrule
\end{tabular}
\end{center}

To run the full pipeline from scratch:
\begin{lstlisting}
doit
\end{lstlisting}

To update only the shock computation and downstream outputs after changing dates:
\begin{lstlisting}
doit compute_yield_shocks make_table_1
\end{lstlisting}

\subsection{Configuration}

All user-facing parameters are stored in \texttt{.env}:

\begin{center}
\begin{tabular}{lll}
\toprule
Variable & Current Value & Description \\
\midrule
\texttt{REPORT\_DATE\_SLASH} & 12/31/2025 & FFIEC reporting quarter \\
\texttt{REPORT\_DATE}        & 12312025     & Date in filename format \\
\texttt{MARKET\_START\_DATE} & 2020-01-01 & Start of shock window \\
\texttt{MARKET\_END\_DATE}   & 2025-12-31   & End of shock window \\
\texttt{RMBS\_MULTIPLIER}    & 1.25   & Convexity haircut for MBS \\
\texttt{FRED\_API\_KEY}      & \textit{(hidden)} & Free key from fred.stlouisfed.org \\
\bottomrule
\end{tabular}
\end{center}

%% ─────────────────────────────────────────────────────────────────────────────
\section{Conclusion}
\label{sec:conclusion}

We replicate the JMPS bank fragility analysis using exclusively open-source data
and a fully automated computational pipeline. Our results confirm the core finding
of JMPS: monetary tightening generates substantial mark-to-market losses across
the U.S. banking system, with the distribution of losses and uninsured deposit
exposure varying significantly across bank size categories.

The primary contribution of this project is methodological. By replacing the
WRDS data dependency with the FFIEC public bulk download, the FRED API, and
Yahoo Finance, we make the analysis accessible to any researcher with a Python
environment and a free FRED API key. The \texttt{doit} pipeline ensures that
updating the analysis to a new Call Report quarter requires only a configuration
change and a single command.

Future extensions could incorporate: (i) panel analysis across multiple Call
Report quarters to track the evolution of fragility over time; (ii) more precise
MBS duration estimates using position-level data from securities databases; and
(iii) integration of interest rate swap positions from Schedule RC-L to account
for hedging activity.

%% ─────────────────────────────────────────────────────────────────────────────
\section*{References}

\begin{itemize}[leftmargin=*, label={}]
    \item Jiang, E., Matvos, G., Piskorski, T., and Seru, A. (2023).
    ``Monetary Tightening and U.S. Bank Fragility in 2023: Mark-to-Market Losses
    and Uninsured Depositor Runs?'' NBER Working Paper 31048.
    \textit{Journal of Political Economy}, forthcoming.

    \item Federal Financial Institutions Examination Council (FFIEC).
    Call Report Bulk Data Download.
    \url{https://cdr.ffiec.gov/public/PWS/DownloadBulkData.aspx}

    \item Federal Reserve Bank of St.\ Louis. FRED Economic Data.
    \url{https://fred.stlouisfed.org}

    \item Diamond, D.\ W.\ and Dybvig, P.\ H.\ (1983).
    ``Bank Runs, Deposit Insurance, and Liquidity.''
    \textit{Journal of Political Economy}, 91(3): 401--419.

    \item Flannery, M.\ J.\ (1981).
    ``Market Interest Rates and Commercial Bank Profitability: An Empirical
    Investigation.''
    \textit{Journal of Finance}, 36(5): 1085--1101.

\end{itemize}

\end{document}
